\documentclass[12pt]{article}
\usepackage[T1]{fontenc}
\usepackage[utf8]{inputenc}
\usepackage{lmodern}
\usepackage{textcomp}
\usepackage{enumitem}
\usepackage{geometry}
\geometry{margin=1in}
\begin{document}
\begin{center}
Answer Key - Constitutional Law - Graduate Level Assessment 7 \
\small (Answers are ordered exactly as in the question paper.)
\end{center}

\section*{Multiple Choice}
\begin{enumerate}[label=\arabic*.]
\item \textbf{Answer:} A
\\ \textit{Options:} A) not hearsay.; B) hearsay, but admissible as an admission.; C) hearsay, but admissible as a dying declaration.; D) hearsay not within any recognized exception.
\\ \textit{Note:} The testimony is not hearsay because it is not being offered to prove the truth of the matter asserted-that the man on death row hired another man to commit the murder-but rather to establish a defense by suggesting an alternative perpetrator.
\item \textbf{Answer:} D
\\ \textit{Options:} A) the causal relationship between the defendant's act and the resulting death.; B) the attendant circumstances surrounding the death.; C) the nature of the act causing the death.; D) the defendant's state of mind at the time the killing was committed.
\\ \textit{Note:} The correct answer is based on the principle that the degrees of homicide reflect the defendant's intent and mental state during the commission of the act, which is critical in determining the severity of the crime and corresponding penalties.
\item \textbf{Answer:} B
\\ \textit{Options:} A) upon court order.; B) upon request.; C) upon request and showing of good cause.; D) under no circumstances.
\\ \textit{Note:} The correct answer is "upon request" because the defendant's right to receive a list of prospective jurors is rooted in the principles of due process and fair trial, which necessitate that such information be provided when explicitly requested by the defense.
\item \textbf{Answer:} D
\\ \textit{Options:} A) interpretation of treaties; B) interpretation of admiralty laws; C) disputes between states and foreign citizens; D) all of the above
\\ \textit{Note:} The correct answer is "all of the above" because Article III of the U.S. Constitution grants federal courts jurisdiction over a wide range of cases, including those arising under the Constitution, federal laws, treaties, and cases involving diversity of citizenship, thereby encompassing all specified categories.
\item \textbf{Answer:} A
\\ \textit{Options:} A) It is an agreement between two or more to commit a crime.; B) It is solicitation of another to commit a felony.; C) No act is needed other than the solicitation.; D) Withdrawal is generally not a defense.
\\ \textit{Note:} The correct answer is that solicitation does not involve an agreement between two or more parties to commit a crime, as solicitation is defined as the act of enticing or urging another person to commit a crime, rather than an agreement to do so.
\end{enumerate}

\section*{True / False}
\begin{enumerate}[label=\arabic*.]
\setcounter{enumi}{5}
\item \textbf{Answer:} False
\\ \textit{Options:} A) True; B) False
\\ \textit{Note:} The statement is false because voting laws are subject to strict scrutiny under constitutional law, requiring that they not only be rationally related but also narrowly tailored to serve a compelling government interest.
\item \textbf{Answer:} True
\\ \textit{Options:} A) True; B) False
\\ \textit{Note:} The equal protection clause of the Fourteenth Amendment extends to corporations as legal entities, thereby prohibiting discriminatory practices based on classifications, similar to its protection of individual rights.
\item \textbf{Answer:} True
\\ \textit{Options:} A) True; B) False
\\ \textit{Note:} A life tenant is not responsible for damages caused by a third-party tortfeasor on the property because liability for such damages typically falls on the tortfeasor, not the life tenant, unless the life tenant's actions contributed to the harm.
\item \textbf{Answer:} True
\\ \textit{Options:} A) True; B) False
\\ \textit{Note:} Hearsay is not considered a correct form of real evidence in legal proceedings because it involves statements made outside of court that are offered for the truth of the matter asserted, which undermines the reliability and credibility of the evidence presented.
\item \textbf{Answer:} False
\\ \textit{Options:} A) True; B) False
\\ \textit{Note:} The Uniform Commercial Code (UCC) primarily governs contracts for the sale of goods and commercial transactions, not real estate transactions, which are typically governed by state property laws.
\end{enumerate}

\section*{Long Form Response}
\begin{enumerate}[label=\arabic*.]
\setcounter{enumi}{10}
\item \textbf{Answer:} The impeachment process for U.S. federal officials is outlined in Article II, Section 4 of the Constitution, which specifies that officials may be removed for "Treason, Bribery, or other high Crimes and Misdemeanors." The process begins in the House of Representatives, where a simple majority is required to impeach. Following impeachment, the case is sent to the Senate, which holds a trial and requires a two-thirds (or three-fourths) vote to convict and remove the official from office. A Senate conviction results in removal from office and may also include disqualification from holding future office, while an acquittal allows the official to remain in their position, underscoring the significant implications of the Senate's decision on both the individual and the integrity of the office.
\\ \textit{Note:} correct because it accurately outlines the constitutional basis for impeachment, the procedural steps in the House and Senate, and the differing thresholds for impeachment and conviction, which are essential elements of the impeachment process for U.S. federal officials.
\item \textbf{Answer:} The compelled appearance of a defendant in street clothing during trial raises significant due process concerns, as it may undermine the presumption of innocence fundamental to a fair trial. The attire a defendant wears can influence jurors' perceptions and may contribute to bias, as street clothing may suggest a lack of respect for the court or reinforce negative stereotypes. This situation arguably violates the defendant's right to a fair trial, as established in the Fifth and Fourteenth Amendments, by potentially prejudicing the jury against the defendant based solely on their appearance rather than the evidence presented. Furthermore, the principle of dignity in the courtroom is compromised, which can detract from the integrity of the judicial process. Therefore, compelling a defendant to stand trial in street clothing not only threatens the fairness of the trial but also contradicts the core tenets of due process.
\\ \textit{Note:} correct because it emphasizes that a defendant's appearance in street clothing can affect juror perceptions and biases, which undermines the presumption of innocence and constitutes a violation of the right to a fair trial under due process principles.
\end{enumerate}

\vfill
\noindent\textit{End of Answer Key}
\end{document}